% !TEX root = ../main.tex
\corrigesChapitre{11}{Structure des molécules organiques}
\corrigesSouples

\begin{blocLivre}{cor-11-1}
\begin{corrige}{1}{Reconnaître les familles fonctionnelles}
\demarche{Repérer le groupe caractéristique : OH, C=O, COOH, COO, NH2 ou CO-NH2.}
\begin{questions}
\item A : groupe hydroxyle, \res{\text{alcool}} (propan-1-ol) ; B : carbonyle, \res{\text{cétone}} (butan-2-one) ;
C : carboxyle, \res{\text{acide carboxylique}} (acide méthanoïque) ; D : \res{\text{ester}} (éthanoate d'éthyle) ;
E : groupe amino, \res{\text{amine}} (éthanamine) ; F : \res{\text{amide}} (éthanamide).
\item Le groupe C$=$O apparaît dans \res{\text{B, C, D et F}} (cétone, acide, ester, amide).
\item D : \res{\ce{C4H8O2}}. L'acide butanoïque \ce{CH3-CH2-CH2-COOH}, acide carboxylique, a la même formule brute :
ester et acide sont ici isomères de constitution.
\end{questions}
\end{corrige}
\end{blocLivre}
\begin{blocLivre}{cor-11-2}
\begin{corrige}{2}{Nommer des molécules}
\demarche{Famille, puis chaîne la plus longue portant le groupe, numérotée pour lui donner le plus petit indice.}
\begin{questions}
\item A : alcool ; B : aldéhyde ; C : cétone ; D : acide carboxylique ; E : halogénoalcane.
\item A : \res{\text{butan-2-ol}} ; B : \res{\text{3-méthylbutanal}} ; C : \res{\text{pentan-2-one}} ;
D : \res{\text{acide 2-méthylpropanoïque}} ; E : \res{\text{1-chlorobutane}}.
\item A : \ce{CH3-CHOH-CH2-CH3} ; D : \ce{(CH3)2CH-COOH}.
\end{questions}
\end{corrige}
\end{blocLivre}
\begin{blocLivre}{cor-11-3}
\begin{corrige}{3}{Du nom à la formule}
\demarche{Le nom code la chaîne, la position du groupe et les ramifications.}
\begin{questions}
\item 2-méthylbutan-1-ol : \ce{HO-CH2-CH(CH3)-CH2-CH3} ; pentan-3-one : \ce{CH3-CH2-CO-CH2-CH3} ; acide
butanoïque : \ce{CH3-CH2-CH2-COOH} ; éthanamide : \ce{CH3-CO-NH2} ; 2-bromopropane : \ce{CH3-CHBr-CH3}.
\item \chemfig[atom sep=1.7em]{-[:30](=[:90]O)-[:-30]O-[:30]-[:-30]}\; : la partie gauche, avec le C$=$O, provient de
l'acide propanoïque ; la partie droite, après l'oxygène, provient de l'éthanol.
\item Pentan-3-one : \res{\ce{C5H10O}} ; 2-méthylbutan-1-ol : \res{\ce{C5H12O}}.
\end{questions}
\end{corrige}
\end{blocLivre}
\begin{blocLivre}{cor-11-4}
\begin{corrige}{4}{Des isomères de constitution}
\demarche{Compter les atomes de chaque squelette, puis comparer les enchaînements.}
\begin{questions}
\item Chaque molécule compte 4 atomes de carbone, 10 d'hydrogène et 1 d'oxygène : même formule brute \ce{C4H10O},
enchaînements différents : ce sont des \res{\text{isomères de constitution}}.
\item A : \res{\text{butan-1-ol}} ; B : \res{\text{butan-2-ol}} ; C : \res{\text{2-méthylpropan-2-ol}}.
\item Dans D, l'oxygène est \res{\text{inséré dans la chaîne}} (pont C--O--C) : la molécule ne possède pas de groupe
hydroxyle, elle n'est pas un alcool.
\item Les isomères diffèrent par exemple par leur température d'ébullition, leur solubilité ou leur réactivité.
\end{questions}
\end{corrige}
\end{blocLivre}
\begin{blocLivre}{cor-11-5}
\begin{corrige}{5}{Lire un spectre infrarouge}
\demarche{Bande large au-dessus de \qty{3200}{cm^{-1}} : O--H ; bandes vers \qty{2900}{cm^{-1}} : C--H.}
\begin{questions}
\item Bande 1, forte et large vers \qty{3350}{cm^{-1}} : liaison \res{\text{O--H}} ; bande 2, vers
\qty{2950}{cm^{-1}} : liaisons \res{\text{C--H}}.
\item La bande O--H signe un \res{\text{alcool}} : avec la formule \ce{C2H6O}, c'est
\res{\text{l'éthanol }\ce{CH3-CH2-OH}}.
\item Non : l'éthanal ne possède pas de O--H, la bande 1 disparaîtrait ; une bande fine et intense apparaîtrait vers
\qty{1730}{cm^{-1}}, celle du C$=$O.
\end{questions}
\end{corrige}
\end{blocLivre}
\begin{blocLivre}{cor-11-6}
\begin{corrige}{6}{Un arôme de poire}
\demarche{Un « …-oate de …-yle » : la première partie vient de l'acide, la seconde de l'alcool.}
\begin{questions}
\item \ce{CH3-COO-CH2-CH2-CH2-CH2-CH3}, soit
\chemfig[atom sep=1.7em]{-[:30](=[:90]O)-[:-30]O-[:30]-[:-30]-[:30]-[:-30]-[:30]}\; : le groupe
\res{\text{ester}} \ce{-COO-}.
\item \res{\ce{C7H14O2}} et $M=7\times12{,}0+14\times1{,}0+2\times16{,}0=\res{\qty{130}{g.mol^{-1}}}$.
\item \ce{CH3-CH2-COO-CH2-CH2-CH2-CH3} : même formule brute, enchaînement différent : ce sont des
\res{\text{isomères de constitution}}.
\item L'\res{\text{acide éthanoïque}} et le \res{\text{pentan-1-ol}}.
\end{questions}
\end{corrige}
\end{blocLivre}
\begin{blocLivre}{cor-11-7}
\begin{corrige}{7}{Trois flacons pour trois spectres}
\demarche{Chercher le C$=$O fin vers \qty{1710}{cm^{-1}}, puis la forme de la bande O--H.}
\begin{questions}
\item Alcool : O--H fort et large (3200--3500) sans C$=$O ; cétone : C$=$O fin et intense (1700--1720) sans O--H ;
acide carboxylique : C$=$O fin et bande O--H très large (2500--3200).
\item Spectre A : C$=$O à \qty{1715}{cm^{-1}}, aucune bande large : \res{\text{propanone}}. Spectre B : bande large à
\qty{3340}{cm^{-1}}, pas de C$=$O : \res{\text{propan-1-ol}}. Spectre C : C$=$O et bande très étalée de 2500 à
\qty{3200}{cm^{-1}} : \res{\text{acide propanoïque}}.
\item Les deux alcools possèdent les mêmes liaisons (O--H, C--H, C--O) : leurs spectres présentent les mêmes bandes,
pratiquement aux mêmes nombres d'onde. L'infrarouge identifie des groupes, pas la longueur de la chaîne.
\end{questions}
\end{corrige}
\end{blocLivre}
\begin{blocLivre}{cor-11-8}
\begin{corrige}{8}{Des groupes dans les médicaments}
\demarche{Chercher OH, C=O, COO, CO–NH sur chaque squelette.}
\begin{questions}
\item Paracétamol : groupe \res{\text{hydroxyle}} \ce{-OH} porté par le cycle et groupe \res{\text{amide}}
\ce{-NH-CO-}.
\item Aspirine : groupe \res{\text{carboxyle}} \ce{-COOH} (acide carboxylique) et groupe \res{\text{ester}}
\ce{-O-CO-}, la famille produite par l'estérification du chapitre~7.
\item Un groupe carboxyle : l'ibuprofène est un \res{\text{acide carboxylique}}, donc un acide faible, qui se
dissocie partiellement dans l'eau selon son $\pKa$.
\item Squelette : $4+1+6+1+1=13$ atomes de carbone ; hydrogènes : $9+2+4+1+1+1=18$ ; deux oxygènes :
\res{\ce{C13H18O2}}.
\end{questions}
\end{corrige}
\end{blocLivre}
\begin{blocLivre}{cor-11-9}
\begin{corrige}{9}{La longueur des chaînes du polyéthylène}
\demarche{La masse d'une chaîne vaut n fois celle du motif.}
\begin{questions}
\item $n\,\ce{CH2=CH2 ->}\ \ce{-(CH2-CH2)_{$n$}-}$ : le motif est \res{\ce{-CH2-CH2-}}.
\item $M(\text{motif})=2\times12{,}0+4\times1{,}0=\res{\qty{28,0}{g.mol^{-1}}}$.
\item Le programme augmente $n$ jusqu'à ce que $n\times28{,}0$ atteigne \num{56000} : il affiche
$\res{n=2000}$ et la masse molaire exacte \qty{56000}{g.mol^{-1}}. Directement :
$n=\dfrac{\num{56000}}{28{,}0}=2000$.
\item Deux carbones par motif : $2\times2000=\res{4000}$ atomes de carbone.
\end{questions}
\end{corrige}
\end{blocLivre}
\begin{blocLivre}{cor-11-10}
\begin{corrige}{10}{Le PVC}
\demarche{Le motif reprend les atomes du monomère, la double liaison en moins.}
\begin{questions}
\item $n\,\ce{CH2=CHCl ->}\ \ce{-(CH2-CHCl)_{$n$}-}$ : motif \res{\ce{-CH2-CHCl-}}.
\item $M(\text{motif})=2\times12{,}0+3\times1{,}0+35{,}5=\res{\qty{62,5}{g.mol^{-1}}}$ ; pourcentage de chlore :
$\dfrac{35{,}5}{62{,}5}=\res{56{,}8~\%}$.
\item $n=\dfrac{\num{1.0e5}}{62{,}5}=\res{1600}$.
\item Le chlorure d'hydrogène est un gaz acide : rejeté, il contribuerait aux pluies acides. Les fumées sont lavées
pour le neutraliser.
\end{questions}
\end{corrige}
\end{blocLivre}
\begin{blocLivre}{cor-11-11}
\begin{corrige}{11}{Un polyamide, le nylon}
\demarche{Le groupe amide relie C=O et N–H ; la masse du motif se calcule sur sa formule brute.}
\begin{questions}
\item Les deux groupes \res{\ce{-CO-NH-}} situés de part et d'autre du motif.
\item $M=12\times12{,}0+22\times1{,}0+2\times14{,}0+2\times16{,}0=\res{\qty{226}{g.mol^{-1}}}$.
\item $n=\dfrac{\num{4.5e4}}{226}\approx\res{200}$ motifs par chaîne.
\item Les groupes amide de chaînes voisines s'attirent (liaisons hydrogène entre N--H et C$=$O) : les chaînes
glissent peu l'une sur l'autre, la fibre est \res{\text{résistante}}.
\end{questions}
\end{corrige}
\end{blocLivre}
\begin{blocLivre}{cor-11-12}
\begin{corrige}{12}{La carte d'identité de l'ibuprofène}
\demarche{Le spectre identifie la famille ; le titrage utilise n = c V à l'équivalence.}
\begin{questions}
\item Bande fine à \qty{1710}{cm^{-1}} : C$=$O ; bande très large de 2500 à \qty{3200}{cm^{-1}} : O--H d'un
\res{\text{acide carboxylique}}, en accord avec le groupe carboxyle de l'ibuprofène.
\item $n=\dfrac{m}{M}=\dfrac{0{,}400}{206{,}3}=\res{\qty{1,94e-3}{mol}}$.
\item À l'équivalence, $n(\ce{HO-})=n(\ce{AH})$ :
$V_\mathrm{E}=\dfrac{n}{c_\mathrm{B}}=\dfrac{\num{1.94e-3}}{0{,}100}=\qty{1,94e-2}{L}=\res{\qty{19,4}{mL}}$.
\item $\pH=6{,}5>\pKa=4{,}9$ : la forme \res{\text{basique }\ce{A-}} prédomine ; le rapport
$[\ce{A-}]/[\ce{AH}]=10^{\,6{,}5-4{,}9}\approx40$.
\end{questions}
\end{corrige}
\end{blocLivre}
\begin{blocLivre}{cor-11-13}
\begin{corrige}{13}{Du PET à la fibre polaire}
\demarche{Quantité de motifs = masse/masse molaire du motif ; degré de polymérisation = M(chaîne)/M(motif).}
\begin{questions}
\item Le motif contient deux groupes \res{\text{ester}} \ce{-O-CO-} : le PET est un polyester.
\item $n(\text{motifs})=\dfrac{25}{192{,}2}=\res{\qty{0,130}{mol}}$.
\item $n=\dfrac{\num{2.4e4}}{192{,}2}=\res{125}$ motifs par chaîne en moyenne.
\item $n(\text{chaînes})=\dfrac{25}{\num{2.4e4}}=\qty{1,0e-3}{mol}$, soit
$\num{1.04e-3}\times\num{6.02e23}\approx\res{\num{6.3e20}}$ chaînes.
\item La refusion dégrade peu à peu les chaînes (elles raccourcissent, le matériau jaunit). Le recyclage chimique
redonne les \res{\text{monomères purs}}, qui permettent de repolymériser un PET de qualité identique au neuf.
\end{questions}
\end{corrige}
\end{blocLivre}
\begin{blocLivre}{cor-11-14}
\begin{corrige}{14}{Les esters et les odeurs de fruits}
\demarche{Dans « …-oate de …-yle », l'acide fournit la partie en -oate et l'alcool la partie en -yle.}
\begin{questions}
\item \chemfig[atom sep=1.7em]{-[:30]-[:-30]-[:30](=[:90]O)-[:-30]O-[:30]-[:-30]}\; : groupe ester \ce{-COO-}.
\item Butanoate d'éthyle : \res{\ce{C6H12O2}} ; éthanoate de butyle \ce{CH3-COO-C4H9} : \res{\ce{C6H12O2}}
également : ce sont des \res{\text{isomères de constitution}}.
\item Acide \res{\text{éthanoïque}} et \res{\text{3-méthylbutan-1-ol}} :
\[\ce{CH3COOH + HO-CH2-CH2-CH(CH3)2 <=> CH3-COO-CH2-CH2-CH(CH3)2 + H2O}.\]
\item $m=\rho V=0{,}87\times5{,}0=\qty{4,4}{g}$ et $M=\qty{130}{g.mol^{-1}}$ :
$n=\dfrac{4{,}35}{130}=\res{\qty{3,3e-2}{mol}}$.
\end{questions}
\end{corrige}
\end{blocLivre}
\begin{blocLivre}{cor-11-15}
\begin{corrige}{15}{Une molécule mystère}
\demarche{La combustion donne les quantités de C et de H ; l'oxygène s'obtient par différence de masse.}
\begin{questions}
\item $n(\ce{C})=n(\ce{CO2})=\dfrac{1{,}76}{44{,}0}=\res{\qty{4,00e-2}{mol}}$ ;
$n(\ce{H})=2\,n(\ce{H2O})=2\times\dfrac{0{,}721}{18{,}0}=\res{\qty{8,0e-2}{mol}}$.
\item $m(\ce{O})=0{,}881-\num{4.00e-2}\times12{,}0-\num{8.0e-2}\times1{,}0=\res{\qty{0,32}{g}}$, soit
$n(\ce{O})=\dfrac{0{,}321}{16{,}0}=\res{\qty{2,0e-2}{mol}}$.
\item Les quantités sont dans les rapports $\ce{C}:\ce{H}:\ce{O}=2:4:1$ : la formule brute est de la forme
$(\ce{C2H4O})_k$, de masse molaire $44{,}0\,k$. Avec $M=\qty{88}{g.mol^{-1}}$, $k=2$ :
$\res{\ce{C4H8O2}}$.
\item C$=$O présent (bande fine à \qty{1740}{cm^{-1}}) mais aucun O--H : ni acide, ni alcool. Parmi les familles en
\ce{C4H8O2}, il reste un \res{\text{ester}}.
\item L'odeur de dissolvant désigne \res{\text{l'éthanoate d'éthyle }\ce{CH3-COO-CH2-CH3}}, compatible avec la
formule brute, la masse molaire et le spectre.
\end{questions}
\end{corrige}
\end{blocLivre}
